% unidades/28-conjunciones-causales.tex
\chapter{➡ \ruby{順接}{じゅんせつ}の\ruby{接続詞}{せつぞくし} — Conjunciones Causales Formales}
\unidadheader{28}{順接の接続詞}{そのため・したがって・その結果}{Expresar causa-efecto en registro formal}

\begin{tcolorbox}[vocabbox, title={\emoji{pushpin} Vocabulario}]
\begin{tabularx}{\linewidth}{llX}
\toprule
\textbf{Japonés} & \textbf{Español} & \textbf{Tipo} \\ \midrule
\vocabitem{結果}{けっか}{resultado / consecuencia}{sustantivo}
\vocabitem{原因}{げんいん}{causa}{sustantivo}
\vocabitem{影響}{えいきょう}{influencia / impacto}{sustantivo}
\vocabitem{導く}{みちびく}{conducir / llevar a}{v. G1}
\vocabitem{証明する}{しょうめいする}{demostrar / probar}{v. suru}
\vocabitem{論理的}{ろんりてき}{lógico}{adj. na}
\vocabitem{解決策}{かいけつさく}{solución}{sustantivo}
\bottomrule
\end{tabularx}
\end{tcolorbox}

\subsection*{\emoji{speech-balloon} Frases Clave}
\frase{\ruby{台風}{たいふう}が\ruby{来}{き}た。そのため、\ruby{学校}{がっこう}は\ruby{休}{やす}みになった。}{Vino un tifón. Por eso, la escuela cerró.}
\frase{\ruby{交通渋滞}{こうつうじゅうたい}があった。したがって、\ruby{遅刻}{ちこく}した。}{Hubo tráfico. En consecuencia, llegué tarde.}
\frase{\ruby{毎日練習}{まいにちれんしゅう}した。その\ruby{結果}{けっか}、\ruby{上達}{じょうたつ}した。}{Practiqué todos los días. Como resultado, mejoré.}
\frase{\ruby{予算}{よさん}が\ruby{足}{た}りない。よって、\ruby{中止}{ちゅうし}する。}{El presupuesto es insuficiente. Por lo tanto, se cancela.}

\subsection*{\emoji{gear} Gramática}
\patron{そのため(に) — «por eso / por esa razón»}
  {Causa. + そのため(に)、+ efecto (neutro/formal)}
  {\ej{\ruby{風邪}{かぜ}をひいた。そのため、\ruby{仕事}{しごと}を\ruby{休}{やす}んだ。}{Me resfrié. Por eso, falté al trabajo.}
   \ej{\ruby{電車}{でんしゃ}が\ruby{止}{と}まった。そのため、\ruby{道}{みち}が\ruby{混}{こ}んだ。}{El tren se detuvo. Por eso, el camino se congestionó.}
   \ej{\ruby{環境}{かんきょう}\ruby{問題}{もんだい}が\ruby{深刻}{しんこく}だ。そのため、\ruby{対策}{たいさく}が\ruby{必要}{ひつよう}だ。}{El problema ambiental es grave. Por eso, son necesarias medidas.}}

\patron{したがって — «por consiguiente / en consecuencia» (muy formal)}
  {Premisa. + したがって、+ conclusión lógica}
  {\ej{この\ruby{実験}{じっけん}は\ruby{失敗}{しっぱい}した。したがって、\ruby{仮説}{かせつ}は\ruby{誤}{あやま}りだ。}{Este experimento fracasó. Por consiguiente, la hipótesis es incorrecta.}
   \ej{\ruby{法律}{ほうりつ}に\ruby{違反}{いはん}している。したがって、\ruby{処罰}{しょばつ}が\ruby{必要}{ひつよう}だ。}{Viola la ley. En consecuencia, se requiere sanción.}
   \ej{\ruby{予算}{よさん}が\ruby{超過}{ちょうか}した。したがって、\ruby{計画}{けいかく}を\ruby{見直}{みなお}す\ruby{必要}{ひつよう}がある。}{El presupuesto se excedió. Por lo tanto, es necesario revisar el plan.}}

\patron{その結果 / よって — resultado concreto / conclusión}
  {Proceso. + その\ruby{結果}{けっか}、+ resultado \quad / \quad よって、+ conclusión (académico/legal)}
  {\ej{\ruby{地道}{じみち}に\ruby{努力}{どりょく}した。その\ruby{結果}{けっか}、\ruby{夢}{ゆめ}を\ruby{実現}{じつげん}できた。}{Me esforcé constantemente. Como resultado, pude realizar mi sueño.}
   \ej{\ruby{研究}{けんきゅう}を\ruby{重}{かさ}ねた。その\ruby{結果}{けっか}、\ruby{新薬}{しんやく}が\ruby{開発}{かいはつ}された。}{Se acumularon investigaciones. Como resultado, se desarrolló un nuevo medicamento.}
   \ej{\ruby{以上}{いじょう}の\ruby{理由}{りゆう}より、よって\ruby{本案}{ほんあん}を\ruby{採択}{さいたく}する。}{Por las razones anteriores, por consiguiente se adopta esta propuesta.}}

\comparacion{そのため vs だから vs ので}
  {そのため / したがって\\\small registro formal / escrito\\\small \ruby{台風}{たいふう}が\ruby{来}{き}た。そのため〜\\\textit{(textos académicos, noticias, informes)}}
  {だから / ので\\\small registro casual / conversación\\\small \ruby{風邪}{かぜ}だから\ruby{休}{やす}む\\\textit{(habla cotidiana, explicaciones directas)}}
  {そのため y したがって conectan oraciones independientes (punto antes) y son propios del registro escrito o formal. だから y ので van dentro de una oración como cláusulas subordinadas y son naturales en conversación. よって es el más formal, usado en textos legales y académicos.}

\coloquial{だから / だもん / だって — causales en habla casual}{
  En conversación diaria las opciones son: だから (neutro), だって (justificación/queja: «es que...»), だもん / だもの (infantil/explicación afectiva). 「行かないの？」「だって、疲れてるんだもん。」 En contextos formales jamás usar estas formas — optar por そのため o したがって.}

\culturabox{\ruby{論理構造}{ろんりこうぞう}と\ruby{日本語}{にほんご} — El argumento formal japonés}{
  En textos académicos y de negocios japoneses, la estructura típica es: premisa → そのため / したがって → conclusión. Esta cadena lógica explícita es el estándar en reportes (レポート), tesis y presentaciones formales. A diferencia del español donde se puede argüir de forma más circular, el japonés formal prefiere cadenas causales lineales y explícitas.
}

\lectura{\ruby{報告書}{ほうこくしょ}のスタイル\\[0.4em]
  「\ruby{今期}{こんき}の\ruby{売上}{うりあげ}は\ruby{予算}{よさん}を\ruby{下回}{したまわ}りました。その\ruby{原因}{げんいん}として、\ruby{市場}{しじょう}\ruby{競争}{きょうそう}の\ruby{激化}{げきか}が\ruby{挙}{あ}げられます。したがって、\ruby{新}{あたら}しい\ruby{戦略}{せんりゃく}が\ruby{必要}{ひつよう}です。\ruby{具体的}{ぐたいてき}な\ruby{対策}{たいさく}を\ruby{実施}{じっし}した\ruby{結果}{けっか}、\ruby{来期}{らいき}には\ruby{改善}{かいぜん}が\ruby{見込}{みこ}まれます。」}
{Estilo de informe\\[0.4em]
  «Las ventas de este período estuvieron por debajo del presupuesto. Como causa, se puede señalar la intensificación de la competencia en el mercado. Por lo tanto, es necesaria una nueva estrategia. Como resultado de implementar medidas concretas, se espera una mejora en el próximo período.»}

\begin{tcolorbox}[ejerciciobox, title={\emoji{pencil} Ejercicios}]
\begin{enumerate}
  \item ¿そのため o したがって? «Llovió mucho. Por eso, hubo inundaciones.» (noticia)
  \item Convierte a formal: 「暑いから、エアコンをつけよう。」
  \item Escribe un párrafo de 3 oraciones usando そのため y その結果.
\end{enumerate}
\textbf{Respuestas:}
1) そのため (\ruby{大雨}{おおあめ}が\ruby{降}{ふ}った。そのため、\ruby{洪水}{こうずい}が\ruby{発生}{はっせい}した。) \quad
2) \ruby{気温}{きおん}が\ruby{高}{たか}い。そのため、エアコンをつける。 \quad
3) (libre)
\end{tcolorbox}

\begin{tcolorbox}[kanjibox, title={\emoji{mahjong} Kanjis}]
\begin{tabularx}{\linewidth}{cllXX}
\toprule
\textbf{Kanji} & \textbf{On} & \textbf{Kun} & \textbf{Significado} & \textbf{Vocab} \\ \midrule
\kanjirow{結}{けつ}{\ruby{むす}{ぶ}}{anudar/resultado}{\ruby{結果}{けっか} / \ruby{結論}{けつろん}}
\kanjirow{果}{か}{\ruby{は}{てる}}{fruto/resultado}{\ruby{結果}{けっか} / \ruby{果物}{くだもの}}
\kanjirow{影}{えい}{\ruby{かげ}{}}{sombra/reflejo}{\ruby{影響}{えいきょう} / \ruby{影}{かげ}}
\kanjirow{響}{きょう}{\ruby{ひび}{く}}{resonar/influir}{\ruby{影響}{えいきょう} / \ruby{反響}{はんきょう}}
\bottomrule
\end{tabularx}
\end{tcolorbox}
